%% Résumé — Ali Alavi (S. A. Alavi Bajestan)
%% Build (from the repo root):
%%   latexmk -pdf -outdir=resume/build resume/Resume-AliAlavi.tex
%%   cp resume/build/Resume-AliAlavi.pdf files/Resume-AliAlavi.pdf
\documentclass[10pt,letterpaper]{article}

\usepackage[T1]{fontenc}
\usepackage[utf8]{inputenc}
\usepackage{charter}
\usepackage[scaled=0.92]{helvet}
\usepackage[letterpaper,top=0.45in,bottom=0.45in,left=0.6in,right=0.6in]{geometry}
\usepackage[protrusion=true,expansion=true]{microtype}
\usepackage{xcolor}
\usepackage{titlesec}
\usepackage{enumitem}
\usepackage{tabularx}
\usepackage{etoolbox}
\usepackage[hidelinks,pdfusetitle]{hyperref}

% single colour: everything is black
\definecolor{accent}{HTML}{000000}
\definecolor{ink}{HTML}{000000}
\definecolor{muted}{HTML}{000000}

\hypersetup{
  pdftitle={Ali Alavi - Resume},
  pdfauthor={Ali Alavi},
  pdfsubject={Resume},
  pdfkeywords={auditory attention decoding, EEG, speech, multimodal, video-language models}
}

\pagestyle{empty}
\setlength{\parindent}{0pt}
\setlength{\tabcolsep}{0pt}
\color{ink}

% ---- keep a heading with at least a few lines of its section (needspace.sty, inlined)
\makeatletter
\newcommand{\needspace}[1]{\par\begingroup\setlength{\dimen@}{#1}%
  \vskip\z@\@plus\dimen@\penalty-100\vskip\z@\@plus-\dimen@\endgroup}
\makeatother

% ---- section headings: sans, bold, hairline rule
\titleformat{\section}
  {\color{accent}\sffamily\bfseries\large}
  {}{0pt}{}[\vspace{-1.5pt}{\titlerule[0.5pt]}]
\titlespacing*{\section}{0pt}{6.5pt}{3pt}
\pretocmd{\section}{\needspace{5\baselineskip}}{}{}

\setlist[itemize]{leftmargin=1.1em,label={\color{accent}\small\textbullet},
  itemsep=1.5pt,parsep=0pt,topsep=1.5pt,partopsep=0pt}
\setlist[enumerate]{leftmargin=1.6em,label={\color{muted}\small[\arabic*]},
  itemsep=1.5pt,parsep=0pt,topsep=1.5pt,partopsep=0pt}

% ---- helpers
% \entry{Organisation}{Location}{Role}{Dates}
\newcommand{\entry}[4]{%
  \begin{tabularx}{\linewidth}{@{}X r@{}}
    \textbf{#1} & {\color{muted}#2} \\
    \textit{#3} & {\color{muted}\small #4}
  \end{tabularx}\par\vspace{1pt}}
% \project{Name}{one-line descriptor}{right-aligned note}
\newcommand{\project}[3]{\textbf{#1} --- #2\hfill{\small #3}}
\newcommand{\me}{\textbf{A.~Alavi}}
\newcommand{\arxiv}[1]{\href{https://arxiv.org/abs/#1}{\color{accent}arXiv:#1}}
\newcommand{\venue}[1]{{\color{muted}\textit{#1}}}

\begin{document}

% =====================================================================
\begin{center}
  {\sffamily\bfseries\fontsize{24}{26}\selectfont Ali Alavi}\\[3pt]
  {\color{muted}\small Ph.D.\ student, Computer Science, The Ohio State University
   \,·\, attention decoding, speech, multimodal \& video-language models}\\[3pt]
  \small
  \href{mailto:alavibajestan.1@osu.edu}{alavibajestan.1@osu.edu} \,|\,
  \href{mailto:sa.alavi.b@gmail.com}{sa.alavi.b@gmail.com} \,|\,
  +1 (812) 360-3048 \,|\,
  \href{https://salialavi.github.io}{\color{accent}salialavi.github.io} \,|\,
  \href{https://github.com/salialavi}{github.com/salialavi} \,|\,
  Columbus, OH
\end{center}
\vspace{-8pt}

% =====================================================================
\section{Education}
\entry{The Ohio State University}{Columbus, OH}{Ph.D.\ in Computer Science}{Sep.\ 2022 -- Present}
{\small Adviser: Dr.\ Donald S.\ Williamson. Research: auditory attention decoding, multimodal EEG and behavioral
sensing, speech quality assessment, neural audio coding.}\par\vspace{3pt}
\entry{Indiana University}{Bloomington, IN}{Ph.D.\ in Computer Science (transferred to OSU)}{Jan.\ 2022 -- Sep.\ 2022}
\vspace{1pt}
\entry{University of Tehran}{Tehran, Iran}{M.Sc.\ in Communication Engineering}{Sep.\ 2016 -- Oct.\ 2020}
{\small Thesis: Monte Carlo simulation of polarization-sensitive optical coherence tomography using parallel programming.}\par\vspace{3pt}
\entry{Ferdowsi University of Mashhad}{Mashhad, Iran}{B.Sc.\ in Electrical and Electronics Engineering}{Sep.\ 2012 -- Mar.\ 2016}
{\small Capstone: LNA compensator with LabVIEW and Xmega microcontrollers.}

% =====================================================================
\section{Research Experience}
\entry{The Ohio State University}{Columbus, OH}{Graduate Research Associate / GTA}{Sep.\ 2022 -- Present}
\begin{itemize}
  \item \project{Merit}{multimodal attention decoding with per-modality credit}{\arxiv{2610.17507}}\\
    Decodes attention from EEG, eye tracking, head motion and scene video; measures what each modality adds.
  \item \project{NEUROTOKEN}{joint source and direction AAD with envelope decoding}{\arxiv{2610.00397}}\\
    One EEG model that identifies the attended talker and direction and reconstructs the envelope via flow matching.
  \item \project{MAESTRO}{open multimodal auditory-attention corpus}{\arxiv{2609.31898}}\\
    Co-developed a public AAD dataset with synchronized EEG, gaze, video and head motion; built the recording setup.
  \item \project{Speech-quality metric hackability}{corrected protocol and benchmark}{\arxiv{2610.10899}}\\
    Measures how easily learned speech-quality metrics can be gamed, and how well common defenses hold up.
  \item \project{CADENZA}{content-adaptive neural audio codec}{\venue{Interspeech 2026 submission}}\\
    Low-bitrate speech codec with variable-rate finite scalar quantization and conditional flow matching.
  \item \project{Contrastive AAD}{self-supervised EEG--speech alignment}{\arxiv{2410.18395}}\\
    Contrastive alignment of EEG and attended speech, fine-tuned to identify the attended talker.
\end{itemize}
\vspace{1pt}
\entry{Indiana University}{Bloomington, IN}{Graduate Research Assistant}{Jan.\ 2022 -- Sep.\ 2022}
\begin{itemize}
  \item Started research on EEG-guided speech enhancement with deep learning under Dr.\ Donald S.\ Williamson.
\end{itemize}

% =====================================================================
\section{Competitions}
\begin{itemize}
  \item \textbf{Silver Medal} --- BirdCLEF+ 2026, Kaggle: species identification from audio recordings.
    \hfill{\small\href{https://www.kaggle.com/certification/competitions/salialavi/birdclef-2026}{\color{accent}certificate}}
  \item \textbf{2nd place, both submitted proposals} --- Obesity ML Challenge Part 3, CrunchDAO, 2026.
    \hfill{\small\href{https://www.linkedin.com/pulse/obesity-ml-challenge-part-3-peer-review-results-crunchlabhq-7slec/}{\color{accent}results}}
  \item \textbf{3rd place} --- CMDC Hackathon.
\end{itemize}

% =====================================================================
\section{Publications \& Preprints}
\begin{enumerate}
  \item \me, D.\,S.\ Williamson. ``Merit: Making every modality earn its contribution in multimodal attention
    decoding.'' \arxiv{2610.17507}, 2026.
  \item \me, D.\,S.\ Williamson. ``How Hackable Is Your Speech Quality Metric? A Corrected Protocol, a Benchmark,
    and What Patching Buys.'' \arxiv{2610.10899}, 2026.
  \item \me, D.\,S.\ Williamson. ``NEUROTOKEN: Joint Source and Directional AAD with Envelope Decoding via
    Conditional Flow Matching.'' \arxiv{2610.00397}, 2026.
  \item K.\,M.\,N.\ Hassan, \me, D.\,S.\ Williamson. ``MAESTRO: a Multimodal Auditory-attention Egocentric
    Speech-TRacking Open corpus.'' \venue{Submitted to IEEE Trans.\ Audio, Speech, and Language Processing};
    \arxiv{2609.31898}, 2026.
  \item \me, D.\,S.\ Williamson. ``CADENZA: Content-Adaptive Encoding with Zero-Collapse Quantization for
    Sub-Kilobit Neural Audio Compression.'' \venue{Submitted to Interspeech}, 2026.
  \item \me. ``TLG: Temporal-Logic Grounding for Video Question Answering via Source-Annotation Reconstruction
    and Category-Targeted Reasoning.'' \venue{TimeLogic Challenge report, CVPR 2026 Workshops}; \arxiv{2606.01591}, 2026.
  \item \me. ``Perception First: A Frontier Native-Video Model with Self-Consistency for Implicit Video Question
    Answering.'' \venue{VRR Challenge report, CVPR 2026 Workshops}; \arxiv{2606.01485}, 2026.
  \item \me. ``Reason, Retrieve, Re-rank: A Zero-Shot Reasoning-Aware Framework for Composed Video Retrieval.''
    \venue{CoVR-R Challenge report, VidLLMs Workshop, CVPR 2026}; \arxiv{2606.00910}, 2026.
  \item \textbf{S.\,A.\ Alavi Bajestan}, M.\ Pitt, D.\,S.\ Williamson. ``A contrastive-learning approach for
    auditory attention detection.'' \arxiv{2410.18395}, 2024.
  \item \me. ``Neural network for estimation of optical characteristics of optically active and turbid
    scattering media.'' \arxiv{2011.06934}, 2020.
\end{enumerate}

% =====================================================================
\section{Industry Experience}
\entry{BNtelligence}{Columbus, OH}{Computer Vision Engineer}{Jul.\ 2025 -- Aug.\ 2025}
\begin{itemize}
  \item Built an end-to-end defect-detection pipeline: data collection and annotation, classification and
    semantic-segmentation models, and a real-time video inference API serving production traffic.
\end{itemize}
\entry{AI-bridge}{Hamburg, Germany}{Deep Learning Consultant}{Oct.\ 2019 -- Mar.\ 2020}
\begin{itemize}
  \item Shipped production deep-learning applications (image enhancement, body-part segmentation,
    super-resolution) in TensorFlow and PyTorch.
\end{itemize}
\entry{Payafanavaran}{Mashhad, Iran}{Software Engineer}{Jan.\ 2019 -- Aug.\ 2019}
\begin{itemize}
  \item Developed a C\# application for 3D surface reconstruction from structured-light scanning data.
\end{itemize}
\entry{Tarashe Sanat-e Pishro}{Mashhad, Iran}{Hardware Design Engineer}{Jul.\ 2016 -- Sep.\ 2016}
\begin{itemize}
  \item Designed a high-speed FPGA board (DDR3 SDRAM, Ethernet, peripherals) for industrial data acquisition.
\end{itemize}

% =====================================================================
\section{Technical Skills}
\begin{tabularx}{\linewidth}{@{}l@{\hspace{8pt}}X@{}}
  \textbf{Research} & auditory attention decoding, EEG/BCI, multimodal fusion and attribution, speech quality
    assessment, neural audio coding, flow matching and diffusion, self-supervised and contrastive learning,
    video-language models, RL for LLMs (GRPO, RLVR) \\[1pt]
  \textbf{Languages} & Python, C/C++, C\#, CUDA, MATLAB, Java, JavaScript \\[1pt]
  \textbf{ML / Infra} & PyTorch, Lightning, TensorFlow, scikit-learn, Hugging Face, Qwen-VL, SigLIP-2, Gemini API,
    Git, Linux, Docker, HPC/Slurm, W\&B, Hydra, \LaTeX, Lean\,4 \\[1pt]
  \textbf{Hardware} & EEG and eye-tracking acquisition and sync, FPGA (VHDL/Verilog), ARM Cortex-M, LabVIEW, Cadence
\end{tabularx}

% =====================================================================
\section{Honors \& Service}
\begin{itemize}
  \item CCBS Summer Research Award, The Ohio State University, 2023
  \item Reviewer, EACL 2021 AfricaNLP Workshop \,·\, Student Volunteer, ICLR 2021, ICML 2021, ICEE 2015
    \,·\, Coordinator, Entrepreneurship Center, Ferdowsi University of Mashhad, 2015--2016
\end{itemize}

\end{document}
